\section*{الفصل \ref{ch:b3:particle-physics} --- فيزياء الجسيمات}

\begin{solution}{exo:b3:particle-physics:1}
(أ) u~u~d: $\tfrac23 + \tfrac23 - \tfrac13 = +1$؛ و u~d~d: $\tfrac23 -
\tfrac13 - \tfrac13 = 0$. (ب) $\tfrac23 + \tfrac13 = +1$ (فنقيض
السفلي يحمل $+\tfrac13$). (ج) $\pi^- =
\bar{\text{u}}\text{d}$؛ و $\Delta^{++} = $ u~u~u ($3\times\tfrac23 =
+2$ --- وهي الحالة التي ساعد إحصاؤها على فرض اختراع
اللون). (د) ومجال الغلوونات الممدود يخزن طاقة
$\propto$ الطول؛ وقبل أن يتحرر \omterm{def:b3:particle-physics:zoo}{كوارك} بزمن طويل، تتجاوز الطاقة المخزونة
$2m_{\text{q}}c^2$ فينقطع الخيط إلى زوج جديد:
فتسحب ميزونًا، لا كواركًا عاريًا قط.
\end{solution}

\begin{solution}{exo:b3:particle-physics:2}
(أ) اللبتونات: e$^-$ و $\nu_\mu$ و $\mu^-$؛ والباريونات: p و n؛ \omterm{def:b3:particle-physics:zoo}{والميزون}:
$\pi^+$. (ب) والهدرونات وحدها: $\pi^+$ و p و n. (ج) والمستقر منفردًا:
e$^-$ و p والنيوترينوات (فالنيوترون الحرّ يعيش خمس عشرة
دقيقة؛ و $\mu$ و $\pi$ يموتان في ميكروثوانٍ وأقلّ). (د) و
المتعادلان: $\nu_\mu$ و n.
\end{solution}

\begin{solution}{exo:b3:particle-physics:3}
(أ) مسموح --- وهو تفكك $\beta^-$ المعياري (يقوم به النيوترون الحرّ،
$T_{1/2} \approx \qty{10}{min}$). (ب) وممنوع بالطاقة:
فالنواتج أثقل من البروتون الحرّ (وداخل نواة تستطيع \omterm{prop:b3:nuclear-physics:binding}{طاقة الارتباط}
أن تدفع، فيحدث تفكك $\beta^+$ المرتبط). (ج)
وممنوع: \omterm{prop:b3:particle-physics:conservation}{فالعدد الباريوني} $1 \to 0$ (\omterm{prop:b3:particle-physics:conservation}{والعدد اللبتوني}
$0 \to -1$). (د) وممنوع \omterm{prop:b3:particle-physics:conservation}{بالعدد اللبتوني} \emph{العائلي}:
إذ كانت الميونية لتختفي وتظهر الإلكترونية --- ولم يُشاهَد
قط، ويُصطاد لهذا السبب بالضبط.
\end{solution}

\begin{solution}{exo:b3:particle-physics:4}
(أ) $6\times0.938 \approx \qty{5.6}{GeV}$. (ب) والطاقة الحركية
$m_{\text{p}}c^2 \approx \qty{0.94}{GeV}$ لكل حزمة
(مع $\sqrt s = 4m_{\text{p}}c^2$ وجهًا لوجه). (ج) ونحو $3:1$:
فحفظ كمية الحركة يحكم على أكثر طاقة الهدف الثابت
بحركة أمامية عديمة النفع للحطام --- أما الحزمتان المتقابلتان فتنفقان
كل شيء. (د) وقد صُمّمت آلة سنة 1955 لتتجاوز عتبة
\qty{5.6}{GeV} بالكاد: فميزانيات المسرّعات تُكتب
بالكتلة اللامتغيرة، ووصل نقيض البروتون في موعده.
\end{solution}

\begin{solution}{exo:b3:particle-physics:5}
(أ) $c\tau = \qty{660}{m}$. (ب) والمدى الممدد البالغ
${\sim}\qty{15}{km}$ يحتاج إلى $\gamma \sim 23$: أي $E = \gamma
m_\mu c^2 \sim \qty{2.4}{GeV}$ --- وهو بالضبط المقياس الكوني. (ج)
وفي مرجعه هو يعيش الميون مدته العادية \qty{2.2}{\micro s}
بينما يمرّ الجوّ مسرعًا مقلَّصًا إلى
${\sim}\qty{650}{m}$: أي بقاء واحد بدفتر آخر. (د) \omterm{prop:b3:particle-physics:conservation}{والعدد
الباريوني}: فلا \omterm{def:b3:particle-physics:zoo}{باريون} خفيف بما يكفي، والميون
\omterm{def:b3:particle-physics:zoo}{لبتون}؛ ويجب أن يغيّر تفككه النكهة، ولا تفعل ذلك إلا القوة الضعيفة
--- ومن هنا الميكروثواني الوقورة بدل يوكتوثواني
القوة القوية.
\end{solution}

\begin{solution}{exo:b3:particle-physics:6}
(أ) $\text{d} \to \text{u} + \text{e}^- + \bar\nu_{\text{e}}$:
فالشحنة $-\tfrac13 = \tfrac23 - 1 + 0$؛ \omterm{prop:b3:particle-physics:conservation}{والعدد الباريوني}
$\tfrac13 = \tfrac13$؛ وعدد العائلة الإلكترونية $0 = 1 - 1$. (ب) و
التفكك الثنائي الجسم يثبّت طاقة الإلكترون؛ فكان \emph{الاتصال}
المشاهَد يعني إما إخفاق حفظ الطاقة
(وقد داعب بور الفكرة) وإما جسمًا ثالثًا غير مرئي يتقاسم
الميزانية --- أي نيوترينو باولي، وقد كُشف بعد خمس وعشرين سنة.
(ج) $\hbar/m_{\text{W}}c = 197/80400 \approx
\qty{2.5e-3}{fm}$ --- أي جزء من ألف من بروتون. (د) وعند الطاقات
المنخفضة يجب على التفاعل أن يجسر فجوةً يجعلها وسيطه الثقيل
قصيرة إلى حدّ العبث، فتكون السعات ضئيلة؛ وقرب
\qty{100}{GeV} يصير W في المتناول وتُظهر القوة ``الضعيفة''
قوّتها الكهرضعيفة الحقيقية.
\end{solution}

\begin{solution}{exo:b3:particle-physics:7}
(أ) القوة القوية تحفظ الغرابة، فلا تستطيع أن تصنع
$+1$ و $-1$ إلا معًا: أي كاون مع هيبرونه. (ب) وثلاث عشرة
رتبة قدر أبطأ من أن تكون للقوة القوية: فالقوة
الضعيفة، وهي كاسرة النكهة الوحيدة، توقّع شهادة الوفاة.
(ج) $10^{13}$: أي ثانية واحدة في مقابل ثلاثمئة ألف سنة.
(د) والولادة الغزيرة قالت ``قوية''؛ والموت المتمنّع قال ``ليست
قوية'' --- فلا يوفّق بين الاثنين إلا عدد محفوظ جديد، تصنعه القوة
القوية أزواجًا ويكسره الضعيف: فمسك الدفاتر
\emph{هو} الاكتشاف.
\end{solution}

\begin{solution}{exo:b3:particle-physics:8}
(أ) في مرجع سكون الزوج تكون كمية الحركة الكلية منعدمة: فلا يستطيع
فوتون واحد أن تكون كمية حركته منعدمة، ويستطيع فوتونان متعاكسان ---
فالحفظ يكتب هندسة التصوير المقطعي البوزتروني. (ب) $E = mc^2 =
0.002\times(3\times10^8)^2 = \qty{1.8e14}{J}$: فغرامان من
الفناء يتفوقان على كيلوغرام يورانيوم مشطور. (ج) ومن
نظائر $\beta^+$ المصنوعة بالسيكلوترون --- أي $^{18}$F، مركَّبًا في
الغلوكوز: فكيمياء الجسم نفسها تضع منبع البوزترون في
أجوع الخلايا. (د) ووقود يكلّف مليون ضعف
مردوده، وبلا مناجم طبيعية: فالمادة المضادّة معايرة
كواشف ومحرك لراوي قصص، لا مصدر طاقة.
\end{solution}

\begin{solution}{exo:b3:particle-physics:9}
(أ) كمون هيغز بئر مزدوجة (أي قبّعة مكسيكية): فالنقطة
المتناظرة $\phi = 0$ قمّة، وقد تدحرج الخلاء
إلى الحافة --- فكسرت الحالة الأساسية للكون تناظرها
هي، تمامًا كما يختار المغنط شماله. (ب) والاختيار
لا يكسر إلا جزءًا من التناظر: فالتركيب الناجي
هو الكهرمغناطيسية، فيحفظ حاملها --- أي الفوتون ---
كتلةً منعدمة، بينما يأكل W و Z، حاملا الاتجاهين المكسورين،
الصلابةَ الموافقة كتلةً. (ج) والاهتزاز القطري
\omterm{def:b3:phase-transitions:order}{لوسيط الترتيب} --- وهو في المغنط تذبذب $|m|$ حول
قيمته التوازنية قرب
$T_{\text{c}}$. (د) وفي أول $10^{-11}$ ثانية من عمر
الكون (\cref{ch:b3:astrophysics}): ففوق درجة الحرارة
الكهرضعيفة كان كل جسيم عديم الكتلة وكانت القوتان
واحدة.
\end{solution}

\begin{solution}{exo:b3:particle-physics:10}
(أ) كل \qty{26.7}{MeV} $= \qty{4.28e-12}{J}$ تعطي
نيوترينوين: فالتدفق $= 2\times1360/4.28\times10^{-12} \approx
\qty{6.4e14}{m^{-2}.s^{-1}}$. (ب) وعبر ظفر:
$6\times10^{10}$ في الثانية. وعبورات العمر كله عبر جسمك
البالغ ${\sim}\qty{0.03}{m^2}$ بسماكة \qty{0.3}{m} على مدى
${\sim}\qty{2.5e9}{s}$: أي نحو $5\times10^{22}$، مضروبةً في
$10^{-22}\,\unit{m^{-1}}\times\qty{0.3}{m}$: أي من رتبة الواحد ---
فنيوترينو شمسي واحد سيلمسك في عمرك كله. (ج)
و $\nu_{\text{e}}$ المفقود كان قد تذبذب إلى $\nu_\mu$ و
$\nu_\tau$ في الطريق؛ والتذبذب يقتضي فروق كتل، ومنه
تكون للنيوترينوات كتلة --- وهو أول صدع مخبري في
النموذج المعياري عديم كتلة النيوترينو. (د) والجبل يرشّح
كلَّ شيء \emph{آخر}: فلا تمرّ إلا النيوترينوات، فتحوّل كيلومترات
الصخر أضجّ سماء إلى أهدأ مخبر.
\end{solution}

\begin{solution}{exo:b3:particle-physics:11}
(أ) $\gamma = 6.8\times10^{12}/9.38\times10^{8} \approx 7250$:
ومنه $1 - v/c \approx 1/2\gamma^2 \approx 10^{-8}$ --- أي ثلاثة أمتار
في الثانية دون الضوء. (ب) $\lambda \approx hc/E \approx
\qty{3e-5}{fm}$: أي جزء من ثلاثين ألفًا من نصف قطر بروتون --- أي أرض
الكواركات. (ج) $13.6\,\text{TeV}/938\,\text{MeV} \approx
14500$ كتلة بروتون --- لو كانت كلها قابلة للتحويل. (د) وعند
طاقة ثابتة تكون $\gamma_{\text{e}}/\gamma_{\text{p}} =
m_{\text{p}}/m_{\text{e}} \approx 1836$، فتكون خسائر الإلكترون
أسوأ بعامل $1836^4 \sim 10^{13}$: فالبروتونات للحلقات ذات القوة الغاشمة، و
الإلكترونات (النظيفة النقطية) للدقة عند طاقة أدنى.
\end{solution}

\begin{solution}{exo:b3:particle-physics:12}
(أ) $3\times10^{20}\times1.6\times10^{-19} \approx \qty{48}{J}$:
أي كرة تنس مضروبة بقوة --- في بروتون واحد. (ب) $\gamma \approx
3\times10^{11}$؛ وسنوات المجرّة الضوئية $10^{5}$ تتقلص إلى
${\sim}10$ \emph{ثوانٍ} ضوئية: فيعبر درب التبانة، بساعته
هو، في ثوانٍ. (ج) وفي مرجع البروتون تصل ميكرويّات
الخلفية الوديعة منزاحةً إلى الأزرق أشعةَ $\gamma$ طاقتها
تكفي لقذف بيونات منه: ففوق العتبة يكون الكون
نفسه ماصًّا، ولا بدّ أن تكون بروتونات كهذه محلية. (د) وجسيم
واحد كهذا لكل كيلومتر مربع لكل قرن، ومن نوع مجهول،
لا يصوّبه أحد: أما تصادمات المصادم الهدروني الكبير البالغة $10^{9}$ في
الثانية، بحزم معلومة عند طاقة معلومة، فتشتري إحصاءً
وتحكمًا لا تضاهيها هدية كونية.
\end{solution}

\begin{solution}{pb:b3:particle-physics:1}
\textbf{1.} بروتون أوّلي (غالبًا \unit{GeV}--\unit{TeV})
يضرب نواة في الأعالي (بالقوة القوية)، فيرشّ بيونات؛ و
البيونات المشحونة تتفكك بالقوة الضعيفة إلى ميونات ونيوترينوات؛ و
الميونات، وهي لا تؤيّن إلا برفق (كهرمغناطيسيًّا)، تنطلق نحو
الأرض.
\textbf{2.} $c\tau_\pi \approx \qty{8}{m}$: فحتى $\gamma$ الكبيرة
لا تشتري للبيونات إلا مئات الأمتار، وهي تموت أيضًا بالتفاعلات
القوية في الهواء؛ أما الميون، وهو لا يحسّ القوة القوية و
يعيش أطول بمقدار $85\times$، فهو الذي يسافر.
\textbf{3.} $400\,\unit{cm^2}\times1\,\text{min}^{-1} = 400$
في الدقيقة --- أي نحو 7 في الثانية.
\textbf{4.} ومضات الضجيج والنشاط الإشعاعي المحلي تشعل لوحًا
واحدًا عشوائيًّا؛ أما اشتعال \emph{كليهما} في حدود \qty{100}{ns}
فيقتضي جسيمًا واحدًا يعبر الزوج --- أي الهندسة
مرشّحًا.
\textbf{5.} $R_{\text{عرضي}} = R_1R_2\Delta t = 100\times100
\times10^{-7} = \qty{e-3}{s^{-1}}$ --- أي عدّة كاذبة واحدة كل
ربع ساعة في مقابل مئة ميون في الدقيقة: أي مهمَلة.
\textbf{6.} لا تُقبل إلا الميونات الواقعة في المخروط الذي يخترق
اللوحين (أي شبه الشاقولية، في حدود ${\sim}30$--$40^\circ$):
فالزوج يقيس اتجاهًا لا مجرد معدّل ---
وذلك ما يجعله \emph{مقرابًا}.
\textbf{7.} والميونات المائلة تعبر جوًّا أكثر: فتفقد طاقةً أكثر
وتقضي زمنًا ذاتيًّا أطول في الطيران، فينجو أقلّ ---
وهو تقريبًا قانون $\cos^2\theta$ الذي ترسمه تجاربك المائلة.
\textbf{8.} $c\tau = \qty{660}{m}$: فالميون الكلاسيكي من
\qty{15}{km} ينجو باحتمال
$\eu^{-15000/660} \approx 10^{-10}$ --- فكانت العليّة لتعدّ
ميونًا واحدًا كل عقد.
\textbf{9.} $\gamma = 4000/106 \approx 38$.
\textbf{10.} والعمر الممدد $\gamma\tau \approx
\qty{84}{\micro s}$: أي مدى ${\sim}\qty{25}{km}$ --- فمعدّل العليّة
هو بالضبط ما تتنبأ به النسبية.
\textbf{11.} وفي مرجع الميون يكون العمر هو
\qty{2.2}{\micro s} العادية لكن الجوّ مقلَّص إلى
$15000/38 \approx \qty{400}{m}$: فيُعبَر ووقت يفيض.
\textbf{12.} $t_{\text{ذاتي}} \approx 15000/(38\times3\times
10^8) = \qty{1.3}{\micro s}$: فالبقاء $\eu^{-1.3/2.2} \approx
0.55$ --- أي يصل النصف، وعدّادك تعدادهم.
\textbf{13.} إنها تجارب الميونات الكلاسيكية بين الجبل ومستوى البحر
(روسي--هول، وفريش--سميث على فيلم): فالتفكك
الكلاسيكي يتنبأ بعدّات منعدمة عمليًّا؛ ومئاتك في
الدقيقة نسبية خاصة، محقَّقةً فوق العزل.
\textbf{14.} \omterm{def:b3:particle-physics:zoo}{لبتون} مشحون من الجيل الثاني:
شحنته $-e$، وسبينه $\tfrac12$، وكتلته \qty{106}{MeV} --- أي
توأم الإلكترون الثقيل.
\textbf{15.} الشحنة: $-1 = -1 + 0 + 0$؛ والعائلة الإلكترونية:
$0 = 1 + (-1) + 0$؛ والعائلة الميونية: $1 = 0 + 0 + 1$. فكل الدفاتر
تتوازن.
\textbf{16.} ثلاثة أجسام تتقاسم الميزانية بكل النسب،
فيكون طيف الإلكترون متصلًا؛ وينتهي حين ترتدّ النيوترينوان
معًا في مواجهة الإلكترون: أي عند نحو
$m_\mu c^2/2 \approx \qty{53}{MeV}$.
\textbf{17.} كان ليكسر \omterm{prop:b3:particle-physics:conservation}{العدد اللبتوني} العائلي --- وهو الدفتر
الوحيد الذي يوازنه النموذج المعياري من غير أن يعرف لماذا. وإيجاده
كان ليكون ابن عمّ تذبذبات النيوترينوات في القطاع المشحون:
أي فيزياء جديدة مضمونة، ولذلك تطارده التجارب حتى
جزء من $10^{13}$.
\textbf{18.} دخول الميون يصنع ومضةً؛ وبعد ميكروثوانٍ
يصنع إلكترون تفككه ومضةً ثانية في اللوح نفسه. و
مدرَّج التأخيرات أُسّيّ ميله \emph{هو}
$\tau = \qty{2.2}{\micro s}$: أي عمر مقيس على رفّ.
\textbf{19.} وكل صيغ بور تتحجّم مع الكتلة: فمدار الهدروجين
الميوني ينكمش بمقدار $207$ إلى ${\sim}\qty{250}{fm}$ ---
أي في عمق سحابة الإلكترونات بحيث يحسّ الميون
حجم البروتون: فالذرات الميونية مساطر للبروتون.
\textbf{20.} $10^{15}/10^{6} = \qty{e9}{eV}$: أي نحو غيغا إلكترون-فولت لكل
جسيم --- فميونات عليّتك مواطنون عاديون في الوابل.
\textbf{21.} والميونات النسبوية تسبق الضوء في الماء فتتوهج
بزرقة تشيرنكوف في مضاعفات الخزان الضوئية. وأثر وابل عند
$10^{20}\,\unit{eV}$ يمتدّ على كيلومترات مربعة
عديدة، فتعاين الخزانات المتفرقة كقطرات مطر في
دلاء.
\textbf{22.} الجسيم يودع طاقةً في وسط
يحوّلها ضوءًا، والمضاعف الضوئي يحوّل الضوء إلى
إلكترونات، والإلكترونيات تعدّها: فكل مسعّر من
العليّة إلى المصادم هو هذه الجملة.
\textbf{23.} نصف قطر الالتفافة يعطي كمية الحركة، وجهتها تعطي
إشارة الشحنة؛ أما الفوتون والنيوترون والنيوترينو ---
وهي متعادلة --- فلا تترك أثرًا، بل ديونها في
مسك الدفاتر لا غير.
\textbf{24.} الطبيعة تسرّع بروتونات مفردة عشرة ملايين
ضعف ما يسرّعه المصادم الهدروني الكبير --- وهو أمر مذلّ --- لكن بمعدل واحد لكل كيلومتر
مربع لكل قرن: فآلاتنا تقايض ذروة الطاقة مقابل
$10^{9}$ تصادمة محكومة في الثانية.
\textbf{25.} مولود من بروتون على ارتفاع عشرة كيلومترات؛ موصَّل
حيًّا \omterm{prop:b3:relativistic-kinematics:dilation}{بتمدد الزمن} ($\gamma \approx 38$، وينجو النصف)؛
معيَّن بومضتين بينهما مئة نانوثانية؛ و
مشروح، كليًّا، بجدول من اثني عشر مقعدًا وأربع
قوى.
\end{solution}
